\chapter{Implementierung}

Dieses Kapitel beantwortet die Umsetzungsfrage der Arbeit: Wie wurde die in Kapitel 4 entworfene Architektur technisch realisiert, so dass der Fahrkern reproduzierbar arbeitet, sicherheitsrelevante Signale Vorrang erhalten und Bedien- und Leitstandsfunktionen den Kernbetrieb nicht stören? Die Darstellung folgt der Roadmap-Struktur und trennt deshalb konsequent zwischen Fahrkern, Sensor- und Sicherheitsbasis, hostseitiger ROS-2-Integration, Sicherheitslogik, Bedien- und Leitstandsebene sowie anschlussfähigen Erweiterungen für intelligente Interaktion.

\section{Hardwareaufbau und elektrische Inbetriebnahme}

Die Hardwareimplementierung verfolgt zwei Ziele. Erstens muss der Fahrkern mechanisch und elektrisch so aufgebaut sein, dass Odometrie und Antriebsregelung reproduzierbare Mess- und Stellgrößen erhalten. Zweitens muss die Sensor- und Sicherheitsbasis so angebunden sein, dass sicherheitsnahe Signale unabhängig von rechenintensiven Host-Funktionen erfasst werden können.

\subsection{Chassis, Antrieb und Grundgeometrie}

Die mechanische Basis bildet ein Differentialantrieb mit zwei koaxial angeordneten Getriebemotoren und einer frei nachlaufenden Stützstruktur. Die präzise Ausrichtung der Motorachsen ist für die Odometrie wesentlich, weil eine fehlerhafte Spurbreite systematische Bahnabweichungen verursacht. Die Antriebsräder wurden deshalb spielfrei auf den Motorwellen montiert. Der verwendete Radradius beträgt

$$r = 32{,}835\,\mathrm{mm}$$

Dieser Wert entspricht dem kalibrierten Raddurchmesser von $65{,}67\,\mathrm{mm}$, der aus Bodentests mit Maßbandvergleich hervorging (siehe Kapitel 6).

Die Stützstruktur wurde so positioniert, dass ein möglichst großer Anteil der Fahrzeugmasse auf den angetriebenen Rädern lastet. Diese Maßnahme reduziert den Einfluss nicht modellierter Roll- und Schlupfeffekte bei Richtungswechseln.

\subsection{Verdrahtung der Dual-Node-Architektur}

Die Low-Level-Ebene wurde mit zwei physisch getrennten XIAO-ESP32-S3-Knoten aufgebaut. Diese Trennung folgt direkt dem Entwurf aus Kapitel 4: Der Drive-Knoten verarbeitet ausschließlich fahrkritische Funktionen, während der Sensor-Knoten I2C-basierte Peripherie und sicherheitsnahe Signale übernimmt. Die Verdrahtung bildet diese Zuständigkeiten eindeutig ab.

Am Drive-Knoten sind die Hall-Encoder mit den Phasen A und B an die Pins D4 und D5 (Phase A) sowie D8 und D9 (Phase B) angeschlossen. Die PWM-Signale für den Cytron-MDD3A-Motortreiber liegen im Dual-PWM-Modus an den Pins D0 bis D3 an. Die Signalleitungen der Encoder wurden räumlich getrennt von den PWM-führenden Motorleitungen verlegt, um kapazitive Einkopplungen der Schaltflanken zu verringern. Die PWM-Frequenz beträgt

$$f_{\mathrm{PWM}} = 20\,\mathrm{kHz}$$

Am Sensor-Knoten bindet der I2C-Bus mit SDA auf D4 und SCL auf D5 die MPU6050, den INA260 und den PCA9685 an. Ein MH-B-Kanten-Sensor ist direkt an einen digitalen GPIO-Pin angeschlossen. Diese Aufteilung vermeidet, dass I2C-Zugriffe den Fahrkern blockieren.

\subsection{Host-nahe Peripherie}

Der Raspberry Pi 5 bindet die hostseitigen Komponenten für Lokalisierung und Kartierung, Navigation, Bedien- und Leitstandsebene sowie spätere Interaktion an. Der RPLIDAR A1 ist per USB angeschlossen und liefert die Distanzdaten für Kartenaufbau und Navigation. Eine frontseitig montierte Sony-IMX296-Kamera ist über CSI angebunden. Für hardwarebeschleunigte Bildverarbeitung wurde ein Hailo-8L-Beschleuniger über ein M.2-HAT integriert. Der Audiopfad nutzt einen MAX98357A-Verstärker als Ausgabestufe.

Diese Peripherie liegt bewusst außerhalb des Low-Level-Regelpfads. Damit bleiben Kamera, Vision und Audio anschlussfähig, ohne die Zykluszeit der Motorregelung direkt zu beeinflussen.

\subsection{Spannungsversorgung und Massekonzept}

Die Spannungsversorgung trennt Rechenlast, Logik und Leistungspfad. Der Raspberry Pi 5 wird über einen Step-Down-Regler mit

$$U = 5\,\mathrm{V}$$

versorgt. Die beiden ESP32-S3-Knoten beziehen ihre Versorgung über USB; die Boards erzeugen die Logikspannung von

$$U = 3{,}3\,\mathrm{V}$$

über interne Regler. Der Motortreiber schaltet die Akkuspannung direkt auf die Antriebe. Der Leistungsbereich liegt damit im Bereich von

$$U = 10{,}8\,\mathrm{V}\ \text{bis}\ 12{,}6\,\mathrm{V}$$

Ein gemeinsamer Sternpunkt für die Masse reduziert Ausgleichsströme und erschwert Masseschleifen zwischen Host, Mikrocontrollern und Leistungsstufe.

\section{Firmware von Fahrkern sowie Sensor- und Sicherheitsbasis}

Die Firmware wurde in zwei getrennten PlatformIO-Projekten umgesetzt. Jedes Projekt läuft auf einem eigenen ESP32-S3-Knoten und nutzt FreeRTOS zur funktionalen und zeitlichen Partitionierung. Diese Trennung überführt die Architekturentscheidung aus Kapitel 4 in lauffähige Firmware.

\subsection{Drive-Knoten für Fahrkern und Odometrie}

Der Drive-Knoten realisiert den Fahrkern. Er verarbeitet Encoder-Signale, berechnet Odometrie und setzt freigegebene Geschwindigkeitsvorgaben in PWM-Signale für den Motortreiber um. Die Software ist in hardwarenahe und regelungstechnische Komponenten gegliedert, darunter \texttt{robot\_hal.hpp}, \texttt{pid\_controller.hpp} und \texttt{diff\_drive\_kinematics.hpp}.

Die Regelaufgabe läuft auf Core 1 als \texttt{controlTask} mit einer Zyklusfrequenz von

$$f_{\mathrm{ctrl}} = 50\,\mathrm{Hz}$$

Die Encoder-Impulse werden per Interrupt erfasst. Der Regler nutzt die Parametrierung

$$K_p = 0{,}4 \qquad K_i = 0{,}1 \qquad K_d = 0{,}0$$

Zur Glättung der Drehzahlmessung verwendet der Knoten einen Exponential-Moving-Average-Filter mit

$$\alpha = 0{,}3$$

Core 0 betreibt den \texttt{rclc\_executor}, publiziert Odometrie auf \texttt{/odom} mit

$$f_{\mathrm{odom}} = 20\,\mathrm{Hz}$$

und empfängt Geschwindigkeitskommandos über \texttt{/cmd\_vel}. Eine gemeinsam genutzte Datenstruktur wird durch einen Mutex geschützt, damit Regelkern und Kommunikationskern konsistente Zustände verarbeiten.

\subsection{Sensor-Knoten für I2C-Peripherie und sicherheitsnahe Signale}

Der Sensor-Knoten realisiert die Sensor- und Sicherheitsbasis. Er liest IMU, Batterieüberwachung und Kanten-Sensorik aus und stellt die Messgrößen dem ROS-2-System über micro-ROS bereit. Die Auslagerung dieser Aufgaben verhindert, dass langsame oder blockierende Sensorzugriffe die Fahrregelung beeinträchtigen.

Die I2C-Abfragen laufen auf Core 1. Da die verwendete Wire-Bibliothek nicht thread-sicher arbeitet, serialisiert ein \texttt{i2c\_mutex} alle Buszugriffe. Der Timeout beträgt

$$t_{\mathrm{mutex}} = 5\,\mathrm{ms}$$

Für die Schätzung der Gierrichtung kombiniert ein Complementary-Filter Gyroskop- und Encoderinformation mit der Gewichtung

$$98\,\%\ \text{Gyroskop} \qquad \text{und} \qquad 2\,\%\ \text{Encoder}$$

Core 0 publiziert IMU-Daten auf \texttt{/imu} mit einer Sollfrequenz von

$$f_{\mathrm{imu,soll}} = 50\,\mathrm{Hz}$$

Die tatsächlich erreichte Frequenz beträgt aufgrund von I2C-Bus-Contention etwa $33\,\mathrm{Hz}$.

Batteriedaten auf \texttt{/battery} mit

$$f_{\mathrm{battery}} = 2\,\mathrm{Hz}$$

und den Kantenstatus auf \texttt{/cliff} mit

$$f_{\mathrm{cliff}} = 20\,\mathrm{Hz}$$

Aktorkommandos, etwa für Servos, werden nicht direkt im Kommunikations-Callback ausgeführt. Stattdessen nutzt der Knoten ein Deferred-Pattern: Der Callback schreibt nur den Sollzustand in den Arbeitsspeicher, die eigentliche I2C-Ausgabe erfolgt in der regulären Ablaufsteuerung. Dieses Muster reduziert Jitter und vermeidet blockierende Buszugriffe im Callback-Kontext.

Der Ultraschallsensor HC-SR04 nutzt eine ISR-basierte, nicht blockierende Messung. Eine Interrupt-Service-Routine auf dem Echo-Pin erfasst die Laufzeit über GPIO-Register-Zugriffe. Das Trigger-Echo-Pattern ersetzt das blockierende \texttt{pulseIn()} und gibt die Rechenzeit des Sensor-Tasks für IMU-Abfragen frei. Die Entfernungsdaten erscheinen auf \texttt{/range/front} mit einer Rate von $10\,\mathrm{Hz}$.

\section{Hostseitige ROS-2-Integration}

Die hostseitige Software läuft auf dem Raspberry Pi 5 mit ROS 2 Humble in einer Container-Umgebung. Diese Ebene bündelt Lokalisierung und Kartierung, Navigation, Benutzeroberfläche, Sicherheitslogik und vorbereitete Interaktionsfunktionen. Die Implementierung hält damit die Trennung zwischen Low-Level-Regelung und rechenintensiver Verarbeitung aus dem Systementwurf ein.

\subsection{micro-ROS-Anbindung der Low-Level-Knoten}

Die beiden ESP32-S3-Knoten sind nicht als proprietäre Peripherie gekoppelt, sondern als micro-ROS-Teilnehmer in den ROS-2-Verbund eingebunden. Zwei getrennte \texttt{micro\_ros\_agent}-Prozesse binden den Drive-Knoten über \texttt{/dev/amr\_drive} und den Sensor-Knoten über \texttt{/dev/amr\_sensor} mit jeweils $921600\,\mathrm{Bd}$ an. Die serielle Kopplung folgt dem in Kapitel 3 begründeten Ziel, den Regelpfad deterministisch und störarm zu halten.

Zusätzlich synchronisiert \texttt{rmw\_uros\_sync\_session(1000)} die Zeitbasis der Mikrocontroller mit der Host-Uhr. Diese Synchronisation ist für konsistente Zeitstempel von Odometrie, Sensorik und Transformationsketten erforderlich.

\subsection{CAN-Bus als redundanter Kommunikationspfad}

Zusätzlich zur UART-basierten micro-ROS-Anbindung übertragen beide ESP32-S3-Knoten ihre Daten parallel über einen CAN-Bus mit $1\,\mathrm{Mbit/s}$ (ISO 11898, SN65HVD230-Transceiver). Die CAN-Sends laufen auf Core 1, damit sie unabhängig vom micro-ROS Agent arbeiten. Der Sensor-Knoten sendet Cliff-Signale (CAN-ID 0x120, 20 Hz) und Unterspannungsereignisse (CAN-ID 0x141) direkt an den Drive-Knoten. Der Drive-Knoten empfängt diese Frames im Regelungs-Task und setzt bei Schutzauslösung die Sollgeschwindigkeiten auf null. Dieser Pfad arbeitet unabhängig vom Raspberry Pi und bildet damit einen hardwarenahen Sicherheitspfad. Auf dem Raspberry Pi empfängt ein optionaler \texttt{can\_bridge\_node} die CAN-Frames via SocketCAN und republiziert sie als ROS-2-Topics.

\subsection{Integration von Lokalisierung und Kartierung sowie Navigation}

Die Implementierung bindet LiDAR, Odometrie und IMU in die hostseitige Verarbeitung ein. Auf dieser Grundlage laufen die Knoten für Lokalisierung und Kartierung sowie Navigation. Der Host verarbeitet damit genau die Funktionsgruppen, die in der Roadmap den Ebenen „Lokalisierung und Kartierung" sowie „Navigation" zugeordnet sind. Die Low-Level-Knoten liefern dafür Zustandsgrößen und sicherheitsnahe Signale, übernehmen jedoch keine globale Zielplanung.

Die Aufgabentrennung hat eine klare Folge: Der Host darf Bewegungsziele, Kartenbezug und Pfadverfolgung berechnen, der Fahrkern setzt ausschließlich freigegebene Sollgrößen um. Damit bleibt die Umsetzungslogik auch bei höherer Systemlast stabil.

\section{Sicherheitslogik und Kommandopfad}

Die Implementierung enthält eine eigene Sicherheitslogik, die reguläre Fahrkommandos überstimmen kann. Diese Logik setzt die in der Roadmap geforderte Trennung zwischen Navigation, Bedienung und Schutzreaktion technisch um.

\subsection{Cliff-Sicherheitsmultiplexer}

Der \texttt{cliff\_safety\_node} bildet den sicherheitsnahen Multiplexer für Kanten- und Hinderniserkennung. Er abonniert \texttt{/cliff} sowie \texttt{/range/front} und überwacht gleichzeitig die eingehenden Fahrkommandos. Erkennt der Sensor-Knoten eine Kante oder unterschreitet die Ultraschall-Distanz $80\,\mathrm{mm}$, unterbricht der Multiplexer den regulären Kommandopfad und sendet stattdessen einen Stopp auf das finale Topic \texttt{/cmd\_vel}. Die Freigabe erfolgt erst bei einer Distanz über $120\,\mathrm{mm}$ (Hysterese).

Der sichere Halt entspricht im Auslösefall der Vorgabe

$$v = 0\,\mathrm{m/s}$$

Diese Umsetzung verschiebt die Kantenerkennung bewusst vor reguläre Navigationsentscheidungen. Eine erkannte Kante ist kein normaler Navigationsfehler, sondern ein Schutzereignis mit Vorrang.

\subsection{Freigegebene Kommandokette}

Die Implementierung folgt nicht dem Muster direkter Motoransteuerung durch beliebige Eingaben. Stattdessen entsteht der Fahrbefehl aus einer freigegebenen Kommandokette. Autonome Navigation, manuelle Bedienung und spätere Sprachbefehle bleiben logisch getrennte Eingangsquellen. Erst nach Freigabe darf ein Ziel- oder Moduskommando in einen Fahrbefehl überführt werden.

Die zulässige Kette lautet damit konzeptionell:

$$\text{Interaktion} \rightarrow \text{Freigabelogik} \rightarrow \text{Missionskommando} \rightarrow \text{Navigation oder Bedienung} \rightarrow /\text{cmd\_vel} \rightarrow \text{Fahrkern}.$$

Diese Struktur verhindert, dass eine spätere Sprachschnittstelle die Sicherheitslogik umgeht. Gleichzeitig bleibt die Bedien- und Leitstandsebene handlungsfähig, weil sie den Betriebszustand beobachten und freigegebene Kommandos auslösen kann.

\section{Bedien- und Leitstandsebene}

Die Bedien- und Leitstandsebene wurde als eigenständiger Implementierungsblock aufgebaut. Sie stellt Telemetrie, Videostream, manuelle Bedienung und Zustandsanzeige bereit und dient damit sowohl dem Betrieb als auch der Diagnose. Diese Ebene greift nicht direkt in die Motorregelung ein, sondern nutzt die definierte Kommandokette.

\subsection{Dashboard-Brücke und Kommunikationspfade}

Das System stellt eine browserbasierte Benutzeroberfläche auf Basis von Vite und React bereit. Ein maßgeschneiderter \texttt{dashboard\_bridge} verbindet ROS-2-Datenpfade mit der Benutzeroberfläche. Im Unterschied zu einer generischen \texttt{rosbridge\_suite} reduziert diese Brücke den Protokoll-Overhead für den konkreten Anwendungsfall.

Die Brücke kombiniert einen WebSocket-Server für Telemetrie auf Port

$$9090$$

mit einem HTTP-Endpunkt für den MJPEG-Stream auf Port

$$8082$$

Die Benutzeroberfläche ist auf dem Host über Port

$$5173$$

erreichbar. Diese Aufteilung trennt zustandsarme Telemetrie von bandbreitenintensivem Videostream.

\subsection{Betriebsfunktionen der Benutzeroberfläche}

Die Benutzeroberfläche visualisiert LiDAR-Daten, Kamerabild, Telemetrie und manuelle Bedienelemente. Damit bildet sie die in der Roadmap geforderte Bedien- und Leitstandsebene mit Diagnosefunktion ab. Die manuelle Steuerung erfolgt über einen Joystick. Zusätzlich stehen Zustandsanzeigen zur Verfügung, die den aktuellen Betriebszustand nachvollziehbar machen.

Ein dreistufiger Deadman-Mechanismus sichert die Teleoperation ab. Die erste Stufe liegt im Browser und sendet beim Loslassen des Joysticks sofort einen Stopp. Die zweite Stufe liegt in der Dashboard-Brücke und stoppt die Ausgabe bei ausbleibendem WebSocket-Heartbeat nach mehr als

$$300\,\mathrm{ms}$$

Die dritte Stufe liegt in der Firmware und greift bei ausbleibenden \texttt{cmd\_vel}-Nachrichten nach mehr als

$$500\,\mathrm{ms}$$

ein. Die drei Stufen wirken nacheinander und erhöhen damit die Robustheit der manuellen Bedienung.

\section{Erweiterungen für Vision, Audio und Sprachschnittstelle}

Die Ebene der intelligenten Interaktion wurde nur teilweise als Kernfunktion umgesetzt. Realisiert wurden die Kameraanbindung, der hardwarebeschleunigte Vision-Pfad und ein Audiopfad. Die Sprachschnittstelle bleibt als vorbereitete Erweiterung anschlussfähig, ist jedoch nicht Teil des fahrkritischen Kernpfads.

\subsection{Hybride Vision-Pipeline}

Die Vision-Verarbeitung wurde aus dem ROS-Echtzeitgraphen ausgelagert, damit Kartierung, Navigation und Sicherheitslogik nicht durch Bildverarbeitung blockiert werden. Der \texttt{v4l2\_camera\_node} erfasst den Kamerastrom. Die Dashboard-Brücke stellt daraus einen MJPEG-Stream bereit. Ein nativer Host-Prozess \texttt{host\_hailo\_runner.py} führt die Objekterkennung auf dem Hailo-8L aus und überträgt erkannte Begrenzungsrahmen per UDP auf Port

$$5005$$

zurück in die Container-Umgebung.

Zusätzlich ist eine asynchrone Semantikstufe mit \texttt{gemini\_semantic\_node} vorgesehen. Diese Stufe bleibt bewusst vom Fahrkern entkoppelt und ergänzt den Wahrnehmungspfad nur oberhalb der Sicherheits- und Freigabelogik.

\subsection{Audioausgabe und vorbereitete Sprachintegration}

Der Audiopfad unterstützt definierte Rückmeldungen an die Bedien- und Leitstandsebene. Die Roadmap ordnet zusätzlich eine Sprachschnittstelle mit ReSpeaker Mic Array v2.0 der Ebene der intelligenten Interaktion zu. Für die Projektarbeit ist entscheidend, dass diese Erweiterung architektonisch vorbereitet bleibt: Sprachbefehle dürfen ausschließlich über Intent-Erkennung, Freigabelogik und Missionskommando auf das System wirken.

Damit gilt auch für die Erweiterung der Grundsatz, dass Sprachverarbeitung keine direkte Motoransteuerung auslösen darf. Die spätere Integrationskette lautet daher:

$$\text{Sprachbefehl} \rightarrow \text{Intent} \rightarrow \text{Freigabelogik} \rightarrow \text{Missionskommando} \rightarrow \text{Navigation, Leitstand oder Audio}.$$

\section{Ergebnis der Implementierung}

Die Implementierung überführt den Entwurf in eine funktionsfähige, klar getrennte Systemstruktur. Zwei ESP32-S3-Knoten realisieren Fahrkern sowie Sensor- und Sicherheitsbasis. Der Raspberry Pi 5 bündelt ROS 2, Lokalisierung und Kartierung, Navigation, Sicherheitslogik und die Bedien- und Leitstandsebene. Kamera, Hailo-8L und Audiopfad erweitern das System oberhalb des fahrkritischen Kerns. Damit liegt eine Implementierung vor, die die Roadmap-Themen in eine konsistente technische Gesamtstruktur überführt und die spätere Validierung gezielt vorbereitet.
